\documentclass[11pt,a4paper]{article}

\usepackage[utf8]{inputenc}
\usepackage[T1]{fontenc}
\usepackage[spanish,es-tabla]{babel}
\usepackage{geometry}
\geometry{margin=2.2cm}
\usepackage{hyperref}
\hypersetup{colorlinks=true,linkcolor=blue!60!black,urlcolor=blue!60!black,pdftitle={Defensa — El Gran Torneo Pokemon}}
\usepackage{xcolor}
\usepackage{amssymb}
\usepackage{enumitem}
\usepackage{underscore}
\usepackage{booktabs}
\usepackage{longtable}
\usepackage{array}
\usepackage{tikz}
\usetikzlibrary{arrows.meta,positioning,calc,shapes.geometric}
\usepackage{listings}
\lstset{basicstyle=\ttfamily\footnotesize,frame=single,language=C,breaklines=true}

\setlength{\parskip}{4pt}

\title{\textbf{El Gran Torneo Pok\'emon}\\[4pt]
\large Gu\'ia de lectura, arquitectura, estructuras de datos,\\ diagramas y verificaci\'on de requisitos para la defensa oral}
\author{Programaci\'on I / Algoritmos I — Universidad de Carabobo\\ C (ISO C99), gcc -std=c99 -Wall -Wextra}
\date{\today}

\begin{document}
\maketitle
\tableofcontents
\newpage

%=====================================================================
\section{C\'omo leer el proyecto: gu\'ia ordenada}
%=====================================================================
El proyecto son 10 m\'odulos (par \texttt{.c}/\texttt{.h}), una cabecera de
constantes y la bater\'ia de pruebas (5\,011 l\'ineas en \texttt{src/}). El orden
recomendado de lectura va de los datos a la l\'ogica, y de lo simple a lo
compuesto; cada paso indica qu\'e buscar.

\begin{enumerate}[leftmargin=1.6em]
\item \textbf{\texttt{src/constantes.h}} (51 l\'ineas) — el ``vocabulario'' del
  dominio: 150 especies, 32 entrenadores, 18 tipos, equipo de 6, niveles 1--100,
  64 combates, buffers y rutas. Sin c\'odigo: solo decisiones en un lugar.
\item \textbf{\texttt{src/pokedex.h/.c}} — la fuente de datos. El struct
  \texttt{Especie}, la carga estricta desde archivo (\texttt{pokedex\_cargar})
  y la inmutabilidad (todo es \texttt{const}).
\item \textbf{\texttt{src/tipos.h/.c}} — el cat\'alogo de 18 tipos y la matriz
  de efectividad $18\times18$ (\texttt{static} en el \texttt{.c}).
\item \textbf{\texttt{src/entrenador.h/.c}} — el struct \texttt{Entrenador}, el
  registro de 32 y el id \'unico.
\item \textbf{\texttt{src/equipo.h/.c}} — el m\'odulo m\'as denso: el
  \texttt{Ejemplar} como nodo de lista enlazada, la f\'ormula de estad\'isticas
  (D2), la separaci\'on especie/ejemplar y el backtracking (\texttt{bt\_rec}).
\item \textbf{\texttt{src/combate.h/.c}} — el motor: f\'ormula de da\~no (D1),
  orden por velocidad, KO, empates y anti-empate; el patr\'on de callback
  (\texttt{CombateSeleccionar}).
\item \textbf{\texttt{src/torneo.h/.c}} — la m\'aquina de estados: grupos A--H,
  calendario 1--48, clasificaci\'on con desempates y bracket 49--64.
\item \textbf{\texttt{src/resultados.h/.c}} — carga por teclado/archivo y el
  cat\'alogo completo de validaci\'on de resultados.
\item \textbf{\texttt{src/archivos.h/.c}} — persistencia: guardar/cargar
  entrenadores, resultados y clasificaci\'on.
\item \textbf{\texttt{src/validacion.h/.c}} — lectura segura con reintentos y
  manejo de EOF (ninguna funci\'on llama \texttt{exit}).
\item \textbf{\texttt{src/main.c}} (794 l\'ineas) — solo flujo y men\'u: 17
  funciones \texttt{static} privadas que orquestan los m\'odulos. Sin l\'ogica
  de negocio.
\item \textbf{\texttt{tests/}} — la evidencia: \texttt{run\_tests.sh}, 20 casos
  de men\'u y 2 probes C (efectividad e integraci\'on completa del torneo).
\end{enumerate}

\emph{Regla de oro del proyecto}: la Pok\'edex jam\'as se escribe; los datos de
las especies son inmutables y los ejemplares copian lo que necesitan.

%=====================================================================
\section{Arquitectura: los m\'odulos}
%=====================================================================
Cada m\'odulo es un TDA: un tipo de dato + sus operaciones, con la cabecera
como contrato p\'ublico y la implementaci\'on encapsulada en el \texttt{.c}.
Los identificadores siguen la convenci\'on \emph{modulo\_accion}.

\begin{longtable}{p{3.2cm} p{4.6cm} p{7.4cm}}
\toprule
\textbf{M\'odulo (.c/.h)} & \textbf{L\'ineas} & \textbf{Responsabilidad y funciones p\'ublicas clave} \\
\midrule
\texttt{constantes.h} & 51 & Dominio en un lugar: \texttt{POKEDEX\_MAX}, \texttt{MAX\_ENTRENADORES}, \texttt{CANT\_TIPOS}, \texttt{MAX\_EQUIPO}, \texttt{NIVEL\_MIN/MAX}, \texttt{TOTAL\_COMBATES}, rutas. \\
\texttt{pokedex} & 308+91 & \texttt{pokedex\_cargar} (void, $+$ok), \texttt{pokedex\_buscar\_numero/nombre} (puras), \texttt{pokedex\_mostrar\_todas/especie}. Posee \texttt{Especie} y \texttt{Pokedex[150]}. \\
\texttt{tipos} & 285+71 & \texttt{tipos\_inicializar}, \texttt{tipos\_es\_valido}, \texttt{tipos\_a\_texto}, \texttt{tipos\_multiplicador} (pura). Posee el enum \texttt{Tipo} y la matriz $18\times18$. \\
\texttt{entrenador} & 80+76 & \texttt{entrenador\_registrar} (void, $+$ok), \texttt{entrenador\_buscar}, \texttt{entrenador\_mostrar\_todos}. Posee \texttt{Entrenador} y \texttt{RegistroEntrenadores[32]}. \\
\texttt{equipo} & 467+195 & \texttt{equipo\_crear\_ejemplar}, \texttt{equipo\_agregar\_ejemplar}, \texttt{equipo\_validar}, \texttt{equipo\_liberar}, \texttt{equipo\_asignar}, \texttt{equipo\_mostrar}, \texttt{equipo\_siguiente\_id}, \texttt{equipo\_formar\_backtracking} (void, $+$ok). Posee \texttt{Ejemplar} (lista enlazada), \texttt{RestriccionesEquipo} y el backtracking. \\
\texttt{combate} & 337+96 & \texttt{combate\_calcular\_danio} (pura), \texttt{combate\_ataca\_primero} (pura), \texttt{combate\_ejecutar} (void, $+$ok). Posee \texttt{ResultadoCombate} y el callback \texttt{CombateSeleccionar}. \\
\texttt{torneo} & 818+219 & \texttt{torneo\_armar\_grupos}, \texttt{torneo\_participantes}, \texttt{torneo\_aplicar\_resultado}, \texttt{torneo\_ordenar\_grupos}, \texttt{torneo\_clasificados}, \texttt{torneo\_clasificacion\_grupo} y 5 \texttt{torneo\_mostrar\_*}. Posee \texttt{Combate[64]} y \texttt{Torneo}. \\
\texttt{resultados} & 398+77 & \texttt{resultados\_cargar\_teclado/archivo} (void, $+$ok), \texttt{resultados\_validar} (void, $+$ok), \texttt{resultados\_mostrar\_pendientes}. \\
\texttt{archivos} & 340+92 & Persistencia: 5 funciones \texttt{archivos\_cargar/guardar\_*} (void, $+$ok). Formatos D8. \\
\texttt{validacion} & 134+82 & \texttt{validar\_leer\_entero/\_msg}, \texttt{validar\_leer\_cadena} (void, $+$ok), \texttt{validar\_separar\_campos} (void, $+$n). Ninguna llama \texttt{exit}. \\
\texttt{main} & 794 & Men\'u de 12 opciones y despacho; 17 ayudantes \texttt{static}; \texttt{main} solo orquesta y libera memoria. \\
\bottomrule
\end{longtable}

\textbf{Convenci\'on de funciones} (especificada en
\texttt{openspec/specs/convencion-funciones/spec.md}, RF-CNV-01..03): toda
funci\'on no-\texttt{void} devuelve \emph{un solo} resultado y no modifica
par\'ametros; las operaciones con m\'ultiples salidas son procedimientos
\texttt{void} con par\'ametros por referencia (patr\'on
\texttt{void f(\dots, bool *exito)}).

%=====================================================================
\section{Estructuras de datos utilizadas}
%=====================================================================

\subsection{Respuesta directa: \textquestiondown pilas? \textquestiondown listas? \textquestiondown colas?}

\begin{itemize}
\item \textbf{Listas: S\'I} — lista enlazada simple de \texttt{Ejemplar} (campo
  \texttt{struct Ejemplar *siguiente}), con inserci\'on al frente y
  \texttt{malloc}/\texttt{free} gestionados por el m\'odulo \texttt{equipo}.
  Es la estructura din\'amica que exige el enunciado.
\item \textbf{Pilas: NO expl\'icitamente} — se usa \emph{la pila del sistema}
  (stack de llamadas) como pila impl\'icita de la recursi\'on del backtracking:
  cada llamada recursiva guarda su estado y el retroceso (backtracking) ``desapila''
  al retornar.
\item \textbf{Colas: NO} — ninguna operaci\'on del torneo requiere orden FIFO.
\end{itemize}

\subsection{Inventario y justificaci\'on}

\begin{longtable}{p{5.4cm} p{3.2cm} p{6.6cm}}
\toprule
\textbf{Estructura} & \textbf{D\'onde} & \textbf{Por qu\'e} \\
\midrule
Arreglo est\'atico \texttt{Pokedex[150]} & pokedex & Tama\~no fijo conocido (150 especies); acceso directo por n\'umero. \\
Matriz \texttt{float[18][18]} de efectividad & tipos (\texttt{static} en el \texttt{.c}) & Acceso $O(1)$ a cada par de tipos; encapsulada tras \texttt{tipos\_multiplicador}. \\
Arreglo \texttt{RegistroEntrenadores[32]} & entrenador & M\'aximo del torneo fijo (D10: exactamente 32). \\
Lista enlazada de \texttt{Ejemplar} & equipo & Tama\~no variable (1..6), se crean/destruyen ejemplares din\'amicamente; demuestra el TDA pedido. \\
\texttt{Combate[64]} & torneo & El torneo tiene exactamente 64 combates numerados (48 grupos + 16 eliminatoria). \\
\texttt{id\_clasificados[16]} & torneo & Los 2 clasificados de cada uno de los 8 grupos (1A, 2A, \dots, 1H, 2H). \\
\texttt{filas[4][7]} (salida) & torneo & Tabla de clasificaci\'on de un grupo para pantalla/archivo. \\
\texttt{especies\_permitidas[150]} & equipo & Restricci\'on de integrantes del backtracking. \\
Pila impl\'icita (recursi\'on) & equipo (\texttt{bt\_rec}) & El backtracking es naturalmente recursivo: el estado parcial vive en el stack de llamadas. \\
\bottomrule
\end{longtable}

%=====================================================================
\section{Diagramas de flujo}
%=====================================================================
Se presentan los cuatro flujos centrales como diagramas y, adem\'as, el flujo
resumido de cada m\'odulo.

\subsection{Carga de la Pok\'edex (\texttt{pokedex\_cargar})}
\begin{center}
\begin{tikzpicture}[node distance=6mm and 4mm,
  box/.style={draw,rounded corners=2pt,align=center,font=\footnotesize,inner sep=4pt},
  dec/.style={draw,diamond,aspect=2,align=center,font=\footnotesize,inner sep=2pt}]
\node[box] (a) {Abrir archivo};
\node[dec,below=of a] (b) {abre?};
\node[box,right=of b,xshift=8mm] (c) {*exito = false; retornar};
\node[box,below=of b] (d) {Leer l\'inea};
\node[dec,below=of d] (e) {\textquestiondown fin de archivo?};
\node[dec,right=of e,xshift=14mm] (f) {\textquestiondown 150 v\'alidas?};
\node[box,right=of f,xshift=10mm] (g) {*exito = true; retornar};
\node[box,below=of g] (k) {cantidad queda 0;\\ *exito=false; retornar};
\node[dec,below=of e] (h) {\textquestiondown 8 campos con `;' y rango v\'alido?};
\node[box,right=of h,xshift=16mm] (i) {Rechazar l\'inea (mensaje con n.\textdegree)};
\node[box,below=of h] (j) {Guardar especie en la tabla};
\draw[-Stealth] (a)--(b);
\draw[-Stealth] (b)--node[above,font=\footnotesize]{no} (c);
\draw[-Stealth] (b)--node[right,font=\footnotesize]{s\'i} (d);
\draw[-Stealth] (d)--(e);
\draw[-Stealth] (e)--node[above,font=\footnotesize]{s\'i} (f);
\draw[-Stealth] (f)--node[above,font=\footnotesize]{s\'i} (g);
\draw[-Stealth] (f.south)|-(k.west);
\draw[-Stealth] (e)--node[right,font=\footnotesize]{no} (h);
\draw[-Stealth] (h)--node[above,font=\footnotesize]{no} (i);
\draw[-Stealth] (i.south)|-(d.east);
\draw[-Stealth] (h)--node[right,font=\footnotesize]{s\'i} (j);
\draw[-Stealth] (j.west)-|(d.west);
\end{tikzpicture}
\end{center}
El archivo mal formado deja la Pok\'edex vac\'ia y el programa contin\'ua
(RF-PDX-06); un combate jam\'as modifica estos datos (RF-PDX-04).

\subsection{Combate (\texttt{combate\_ejecutar})}
\begin{center}
\begin{tikzpicture}[node distance=5mm and 4mm,
  box/.style={draw,rounded corners=2pt,align=center,font=\footnotesize,inner sep=4pt},
  dec/.style={draw,diamond,aspect=2.4,align=center,font=\footnotesize,inner sep=2pt}]
\node[box] (a) {Restaurar hp de ambos equipos};
\node[box,below=of a] (b) {Seleccionar activo de cada uno\\(callback: ``local'' / ``visita'')};
\node[dec,below=of b] (c) {EOF o sin equipo?};
\node[box,right=of c,xshift=12mm] (d) {*exito=false (cancelado)};
\node[dec,below=of c] (e) {\textquestiondown turno 20 alcanzado?};
\node[dec,below=of e,yshift=-2mm] (f) {Ataca el m\'as veloz (D3);\textquestiondown cae a 0?};
\node[box,right=of f,xshift=14mm] (g) {Sumar KO; pedir reemplazo\\(``reemplazo'')};
\node[dec,below=of f,yshift=-2mm] (h) {\textquestiondown alguien sin disponibles?};
\node[box,below=of h] (i) {Ganador o empate (D4)\\o anti-empate (D5): HP total $\to$ nivel $\to$ local};
\node[box,right=of i,xshift=8mm] (j) {Escribir *res y *exito=true};
\draw[-Stealth] (a)--(b);
\draw[-Stealth] (b)--(c);
\draw[-Stealth] (c)--node[above,font=\footnotesize]{s\'i}(d);
\draw[-Stealth] (c)--node[right,font=\footnotesize]{no}(e);
\draw[-Stealth] (e)--node[right,font=\footnotesize]{no}(f);
\draw[-Stealth] (f)--node[above,font=\footnotesize]{s\'i}(g);
\draw[-Stealth] (g.south)|-(f.east);
\draw[-Stealth] (f)--node[right,font=\footnotesize]{no}(h);
\draw[-Stealth] (h)--node[right,font=\footnotesize]{s\'i}(i);
\draw[-Stealth] (h.west)-|node[left,font=\footnotesize]{no}(e.west);
\draw[-Stealth] (i)--(j);
\end{tikzpicture}
\end{center}
El da\~no usa la f\'ormula D1 y el multiplicador de tipos
($\times2/\times1/\times0.5/\times0$, producto si el defensor tiene dos tipos).

\subsection{Backtracking de formaci\'on de equipos (\texttt{bt\_rec})}
\begin{center}
\begin{tikzpicture}[node distance=5mm and 4mm,
  box/.style={draw,rounded corners=2pt,align=center,font=\footnotesize,inner sep=4pt},
  dec/.style={draw,diamond,aspect=2.6,align=center,font=\footnotesize,inner sep=2pt}]
\node[box] (a) {bt\_rec(i, restantes, acumuladores, parcial)};
\node[dec,below=of a] (b) {restantes = 0?};
\node[dec,below=of b] (c) {\textquestiondown cumple tipos y stats?};
\node[box,right=of c,xshift=14mm] (d) {*exito=true (primera soluci\'on)};
\node[box,left=of c,xshift=-14mm] (e) {retroceder: *exito=false};
\node[dec,below=of c] (f) {\textquestiondown alguna poda corta?\\(cantidad, nivel, tipos, integrantes)};
\node[box,right=of f,xshift=16mm] (g) {retroceder sin explorar};
\node[box,below=of f] (h) {Incluir k copias de la especie i\\(crear ejemplares con nivel factible)};
\node[box,below=of h] (i) {bt\_rec(i+1, restantes-k, \dots)};
\node[dec,below=of i] (j) {*exito?};
\node[box,left=of j,xshift=-16mm] (k) {Liberar las k copias (sin fugas);\\excluir especie i: bt\_rec(i+1, \dots)};
\draw[-Stealth] (a)--(b);
\draw[-Stealth] (b)--node[right,font=\footnotesize]{s\'i}(c);
\draw[-Stealth] (c)--node[above,font=\footnotesize]{s\'i}(d);
\draw[-Stealth] (c)--node[above,font=\footnotesize]{no}(e);
\draw[-Stealth] (b)--node[right,font=\footnotesize]{no}(f);
\draw[-Stealth] (f)--node[above,font=\footnotesize]{s\'i}(g);
\draw[-Stealth] (f)--node[right,font=\footnotesize]{no}(h);
\draw[-Stealth] (h)--(i);
\draw[-Stealth] (i)--(j);
\draw[-Stealth] (j)--node[above,font=\footnotesize]{no}(k);
\draw[-Stealth] (k.north)|-(f.east);
\end{tikzpicture}
\end{center}
Las cinco podas (cantidad, cota inferior de nivel, cota de tipos con sufijo,
tipos por ejemplar restante y restricci\'on de integrantes) hacen que el caso
imposible (6 ejemplares, nivel total m\'aximo 5) se descarte en microsegundos.
Sin soluci\'on, se informa y \emph{no} se devuelve equipo inv\'alido (RF-EQP-04).

\subsection{M\'aquina de estados del torneo}
\begin{center}
\begin{tikzpicture}[node distance=8mm and 6mm,
  box/.style={draw,rounded corners=2pt,align=center,font=\footnotesize,inner sep=4pt}]
\node[box] (a) {SIN\_INICIAR};
\node[box,right=of a,xshift=6mm] (b) {GRUPOS\\(calendario 1--48)};
\node[box,right=of b,xshift=6mm] (c) {ELIMINATORIAS\\(bracket 49--64)};
\node[box,right=of c,xshift=6mm] (d) {FINALIZADO\\(campe\'on G64)};
\draw[-Stealth] (a)--node[above,font=\footnotesize]{armar\_grupos (32 exactos)} (b);
\draw[-Stealth] (b)--node[above,font=\footnotesize]{48/48 aplicados $\to$ clasificados} (c);
\draw[-Stealth] (c)--node[above,font=\footnotesize]{64 aplicado $\to$ posiciones finales} (d);
\end{tikzpicture}
\end{center}
Los participantes de cada combate los resuelve el sistema (RF-RES-03); el
usuario nunca los inventa: 49 = 1A-2B, \dots, 57 = G49-G50, \dots,
63 = P61-P62, 64 = G61-G62.

\subsection{Flujo resumido de cada m\'odulo}
\begin{longtable}{p{3.0cm} p{12.2cm}}
\toprule
\textbf{M\'odulo} & \textbf{Flujo principal} \\
\midrule
\texttt{pokedex} & cargar (archivo $\to$ parseo estricto $\to$ tabla) $\to$ consultas de solo lectura (mostrar, buscar por n\'umero/nombre, ficha). \\
\texttt{tipos} & inicializar (cargar matriz o respaldo embebido) $\to$ validar nombres $\to$ consultar multiplicador. \\
\texttt{entrenador} & registrar (validar id $\to$ rechazar duplicado $\to$ insertar) / buscar / listar. \\
\texttt{equipo} & crear ejemplar (especie $\to$ f\'ormula D2) $\to$ agregar (lista) $\to$ validar; o formar autom\'aticamente (restricciones $\to$ backtracking $\to$ asignar). \\
\texttt{combate} & restaurar $\to$ seleccionar $\to$ turnos (da\~no D1, KO, reemplazo) $\to$ fin (empate/anti-empate) $\to$ escribir resultado. \\
\texttt{torneo} & armar grupos $\to$ calendario $\to$ aplicar resultados $\to$ ordenar/desempates $\to$ clasificados $\to$ bracket $\to$ campe\'on. \\
\texttt{resultados} & leer (teclado/archivo) $\to$ validar contra el estado real $\to$ aplicar o rechazar con motivo. \\
\texttt{archivos} & guardar/cargar entrenadores, resultados y clasificaci\'on con los formatos D8 (separador `;'). \\
\texttt{validacion} & leer $\to$ reintentar ante entrada inv\'alida $\to$ devolver valor o EOF (cierre ordenado). \\
\texttt{main} & cargar datos $\to$ men\'u $\to$ despachar a m\'odulos $\to$ guardar si hubo cambios $\to$ liberar memoria $\to$ salir. \\
\bottomrule
\end{longtable}

%=====================================================================
\section{Modelo entidad-relaci\'on}
%=====================================================================

\begin{center}
\begin{tikzpicture}[node distance=12mm and 8mm,
  ent/.style={draw,rectangle,rounded corners=2pt,align=center,font=\footnotesize,inner sep=4pt},
  rel/.style={draw,diamond,aspect=2,align=center,font=\footnotesize,inner sep=2pt,fill=blue!5}]
\node[ent] (esp) {ESPECIE\\ \footnotesize numero (PK), nombre,\\ tipo1, tipo2, stats base};
\node[ent,below=of esp,yshift=-4mm] (ej) {EJEMPLAR\\ \footnotesize id (PK), apodo, nivel,\\ stats D2, hp\_actual,\\ tipo1, tipo2, *siguiente};
\node[ent,right=of ej,xshift=10mm] (entr) {ENTRENADOR\\ \footnotesize id (PK), nombre,\\ V/E/D, puntos,\\ pokemon\_derrotados,\\ *equipo};
\node[ent,right=of entr,xshift=10mm] (torn) {TORNEO\\ \footnotesize estado,\\ combates[64],\\ id\_clasificados[16]};
\node[ent,right=of torn,xshift=10mm] (comb) {COMBATE\\ \footnotesize numero (PK 1..64),\\ ent1, ent2, ganador,\\ empate, kos1, kos2};
\node[rel,below=of ej,yshift=-4mm] (r1) {DERIVA\\ 1:N};
\node[rel,below=of entr,yshift=-4mm] (r2) {POSEE\\ 1:N};
\node[rel,below=of torn,yshift=-4mm] (r3) {CONTIENE\\ 1:64};
\node[rel,below=of comb,yshift=-4mm] (r4) {PARTICIPA\\ 2:N};
\draw[-Stealth] (esp)--node[right,font=\footnotesize]{origen de} (ej);
\draw[-Stealth] (entr)--node[left,font=\footnotesize]{due\~no de} (ej);
\draw[-Stealth] (torn)--node[left,font=\footnotesize]{compone} (comb);
\draw[-Stealth] (entr)--node[above,font=\footnotesize]{juega} (comb);
\end{tikzpicture}
\end{center}

\begin{longtable}{p{3.6cm} p{5.6cm} p{6.0cm}}
\toprule
\textbf{Relaci\'on} & \textbf{Cardinalidad} & \textbf{Implementaci\'on en C} \\
\midrule
ESPECIE--EJEMPLAR & 1:N (derivaci\'on) & El ejemplar referencia la especie por \texttt{numero\_especie} y \emph{copia} nombre, tipos y stats (f\'ormula D2) — la especie jam\'as se modifica. \\
ENTRENADOR--EJEMPLAR & 1:N (composici\'on) & Lista enlazada: \texttt{Entrenador.equipo} apunta a la cabeza; cada \texttt{Ejemplar.siguiente} al pr\'oximo. \\
TORNEO--COMBATE & 1:64 (composici\'on) & Arreglo \texttt{Combate[64]} dentro de \texttt{Torneo}; numeraci\'on fija 1--48 grupos, 49--64 eliminatoria. \\
COMBATE--ENTRENADOR & 2:N (participaci\'on) & Cada combate guarda \texttt{id\_entrenador1/2} resueltos por el sistema, ganador y KOs. \\
RESULTADO--COMBATE & 1:1 (aplicaci\'on) & \texttt{ResultadoCargado} valida y aplica sobre el combate correspondiente; parciales permitidos. \\
ARCHIVOS & persistencia externa & \texttt{data/*.txt}: pok\'edex (solo lectura), entrenadores, resultados y clasificaci\'on (salida). \\
\bottomrule
\end{longtable}

%=====================================================================
\section{Cumplimiento del documento de requerimientos}
%=====================================================================
Verificaci\'on contra \emph{Proyecto- 2026.pdf}, secci\'on por secci\'on.
Evidencia ejecutable: \texttt{bash tests/run\_tests.sh} (23 PASS / 0 FALLA),
\texttt{bash scripts/verificar\_entrega.sh} (7 PASS) y build con
\texttt{gcc -std=c99 -Wall -Wextra} con cero warnings.

\begin{longtable}{p{5.2cm} p{2.0cm} p{8.0cm}}
\toprule
\textbf{Requisito (secci\'on)} & \textbf{Estado} & \textbf{Evidencia} \\
\midrule
1.1 Pok\'edex 150 especies, 18 tipos, inmutable & $\checkmark$ & \texttt{data/pokedex.txt} (150 l\'ineas, 0 duplicados); API \texttt{const}; sha256 intacto tras consultas y combates. \\
1.2 Consulta de la Pok\'edex & $\checkmark$ & Opci\'on 1: mostrar todas, buscar por n\'umero y por nombre (insensible a may\'usculas/acentos), ficha completa. \\
1.3 Entrenadores (id \'unico, V/E/D, puntos) & $\checkmark$ & Opci\'on 2 con rechazo de duplicados; struct completo. \\
1.4 Equipos y ejemplares (separaci\'on especie/ejemplar) & $\checkmark$ & \texttt{equipo\_crear\_ejemplar} (f\'ormula D2 con variaci\'on por id); la especie no se muta. \\
1.5 Backtracking & $\checkmark$ & Opci\'on 3 autom\'atica; 5 podas; caso sin soluci\'on informa sin equipo inv\'alido. \\
1.6 Combates & $\checkmark$ & Motor completo: velocidad, da\~no, KO, reemplazo, agotamiento. \\
1.7 Efectividad de tipos & $\checkmark$ & Matriz 18$\times$18; valores $\times2/\times1/\times0.5/\times0$ y producto para dos tipos. \\
1.8 Fase de grupos & $\checkmark$ & 8 grupos de 4 (32 exactos), round-robin 1--48 sin repetidos. \\
1.9 Puntuaci\'on y desempates & $\checkmark$ & 3/1/0; cadena puntos $\to$ victorias $\to$ KOs $\to$ directo (D9) $\to$ id. \\
1.10 Clasificaci\'on autom\'atica & $\checkmark$ & 1A/2A\dots1H/2H resueltos por el sistema al completar 48/48. \\
1.11 Eliminatorias 49--64 & $\checkmark$ & Bracket fijo; G\#/P\# encadenados; sin empates (anti-empate D5); campe\'on G64; top 4. \\
1.12/1.16 Resultados (teclado/archivo, parciales) & $\checkmark$ & Opci\'on 6; \texttt{data/resultados.txt} parcial recargable. \\
1.13 Validaci\'on de resultados & $\checkmark$ & Cat\'alogo completo: existan, pertenezcan, ronda correcta, clasificados, eliminados no reaparecen, ganador v\'alido, sin empates, coherencia entre rondas. \\
1.14 Archivo de la Pok\'edex & $\checkmark$ & 150 especies le\'idas del archivo; cero datos de especies en el c\'odigo. \\
1.15/1.17 Archivos de entrenadores/clasificaci\'on & $\checkmark$ & Guardar/cargar con equipos reconstruibles; \texttt{clasificacion.txt} con formato exacto. \\
2. Men\'u principal & $\checkmark$* & Las 12 opciones visibles; \emph{ver observaci\'on}. \\
3. Implementaci\'on (funciones, structs, arreglos/din\'amicas, archivos, cadenas, TDA, recursividad, backtracking, m\'odulos, main flujo) & $\checkmark$ & 10 m\'odulos; lista enlazada; \texttt{bt\_rec} recursivo; main sin l\'ogica de negocio. \\
4. Validaci\'on de entradas & $\checkmark$ & \texttt{validacion.c}; ninguna entrada termina el programa inesperadamente (bater\'ia E2E). \\
6. Pruebas del sistema (16 puntos) & $\checkmark$ & Cobertura 16/16 en \texttt{tests/} (ver trazabilidad). \\
7. Entrega (un comprimido, compila/ejecuta sin externos) & $\checkmark$ & \texttt{make entrega} $\to$ zip de 93 archivos; \texttt{verificar\_entrega.sh}: 7 PASS en directorio limpio. \\
8. Evaluaci\'on (13 criterios + defensa) & $\checkmark$ & Cada criterio mapeado en \texttt{docs/trazabilidad-requisitos.md}; este documento prepara la defensa. \\
\bottomrule
\end{longtable}

\textbf{Observaci\'on honesta (la \'unica)}: la opci\'on 9 del men\'u
(``Consultar historial de combates'') se muestra pero responde
``en construcci\'on''. El enunciado (secci\'on 2) permite modificar el men\'u;
se recomienda, antes de la entrega, implementarla (registro de combates de la
sesi\'on) o retirarla del men\'u para no dejar funciones sin terminar.

%=====================================================================
\section{Guion de defensa: preguntas probables}
%=====================================================================
\begin{enumerate}[leftmargin=1.6em]
\item \textbf{\textquestiondown Por qu\'e tantas funciones \texttt{static} en \texttt{main.c}?}
  Enlace interno: privadas del archivo, evitan colisiones de nombres y
  encapsulan (es el \texttt{private} de C). Se demuestra con \texttt{nm}: los
  \texttt{static} son s\'imbolos locales \texttt{t} y los p\'ublicos \texttt{T}.
\item \textbf{\textquestiondown Por qu\'e la Pok\'edex nunca se modifica?}
  Separaci\'on especie/ejemplar (criterio central): el ejemplar copia lo que
  necesita y solo muta \texttt{hp\_actual}; las funciones de consulta son
  \texttt{const}-correctas y el sha256 del archivo queda intacto.
\item \textbf{\textquestiondown De d\'onde salen las estad\'isticas de un ejemplar?}
  F\'ormula propia documentada (D2): \texttt{(base*nivel/50)+nivel+5/10}
  m\'as variaci\'on determinista \texttt{(id*7)\%16} en [0,15]: variabilidad
  entre ejemplares y pruebas 100\% reproducibles (libertad otorgada por el
  docente).
\item \textbf{\textquestiondown Por qu\'e no usaste la f\'ormula oficial de da\~no?}
  La oficial incluye IVs/EVs y el t\'ermino del poder del movimiento (fuera de
  alcance); sin \'el, los combates ser\'ian interminables. La f\'ormula D1 usa
  el m\'inimo exigido (nivel, ataque, defensa, tipos) y est\'a justificada en
  el informe.
\item \textbf{\textquestiondown C\'omo evita el backtracking explotar con 150 especies?}
  Cinco podas (cantidad, cota inferior de nivel, cota de tipos, tipos por
  ejemplar restante e integrantes permitidos); el caso imposible se descarta en
  microsegundos y no devuelve equipo inv\'alido.
\item \textbf{\textquestiondown Qu\'e estructuras de datos usaste y por qu\'e?}
  Arreglos est\'aticos donde el tama\~no es fijo (Pok\'edex 150, 32
  entrenadores, 64 combates), matriz 18$\times$18 para efectividad y lista
  enlazada para los equipos (tama\~no variable). No hay pila ni cola
  expl\'icitas: la recursi\'on del backtracking usa la pila del sistema.
\item \textbf{\textquestiondown C\'omo validas que un resultado del torneo es correcto?}
  Ocho chequeos (existen, pertenecen al combate, ronda correcta, clasificados,
  eliminados, ganador v\'alido, sin empates en eliminatoria, coherencia entre
  rondas) y los participantes siempre los resuelve el sistema.
\item \textbf{\textquestiondown Por qu\'e \texttt{main} no tiene l\'ogica de negocio?}
  Cada m\'odulo es un TDA; \texttt{main} solo despacha. El combate ni siquiera
  lee la consola: recibe un callback (\texttt{CombateSeleccionar}) que
  \texttt{main} provee — por eso los probes pueden probar combates sin teclado.
\item \textbf{\textquestiondown Qu\'e convenci\'on siguen tus funciones?}
  RF-CNV-01..03: cada funci\'on devuelve un solo resultado; las operaciones con
  varias salidas son procedimientos \texttt{void} con par\'ametros por
  referencia (p. ej. \texttt{void pokedex\_cargar(\dots, bool *exito)}).
\end{enumerate}

%=====================================================================
\section{Comandos \'utiles durante la defensa}
%=====================================================================
\begin{itemize}[leftmargin=1.6em]
\item \texttt{make} / \texttt{make run} / \texttt{make clean} — compilar, ejecutar, limpiar.
\item \texttt{bash tests/run\_tests.sh} — bater\'ia completa (23 PASS / 0 FALLA).
\item \texttt{bash scripts/verificar\_entrega.sh} — compilar y ejecutar el zip en un directorio limpio.
\item \texttt{make entrega} — regenerar el paquete \texttt{dist/gran-torneo-pokemon-entrega.zip}.
\item \texttt{nm build/torneo | grep -E " [tT] "} — s\'imbolos locales (\texttt{t}, static) vs globales (\texttt{T}).
\item \texttt{doxygen Doxyfile} — documentaci\'on HTML del c\'odigo (binario local en \texttt{tools/} o \texttt{sudo apt install doxygen}).
\item \texttt{grep -h "\#\#\# Requirement" openspec/specs/*/spec.md} — los 45 requisitos formales con su c\'odigo.
\end{itemize}

\vspace{6mm}
\noindent\rule{\textwidth}{0.4pt}
\begin{center}\small
Documento generado como material de estudio para la defensa oral del proyecto
\textbf{El Gran Torneo Pok\'emon} (Programaci\'on I, Universidad de Carabobo).
\end{center}

\end{document}
